\section{Элементы комбинаторики}

\begin{definition}{Размещение}{}
    Пусть дано множество из $n$ элементов.
    Каждое его \textit{упорядоченное} подмножество из $k$ элементов
    называется \textbf{размещением} из $n$ по $k$ и обозначается $\Ank{n}{k}$.
\end{definition}

Количество размещений:
\[
    \Ank{n}{k} = n(n-1)(n-2) \cdots (n-(k-1)) = \frac{n!}{(n-k)!}
\]

Факториал числа $n$:
\[
    n! \coloneqq 1 \cdot 2 \cdot 3 \cdots n, \text{ причем } 0! \coloneqq 1
\]

\begin{definition}{Перестановка}{}
    \textbf{Перестановкой} называется размещение из $n$ по $n$ и обозначается $\Pn{n}$.
\end{definition}

Количество перестановок:
\[
    \Pn{n} = n!
\]

\begin{definition}{Сочетание}{}
    Пусть дано множество из $n$ элементов.
    Каждое его \textit{неупорядоченное} подмножество из $k$ элементов
    называется \textbf{сочетанием} из $n$ по $k$ и обозначается $\Cnk{n}{k}$.
\end{definition}

Количество сочетаний:
\[
    \Cnk{n}{k} = \frac{n(n-1)(n-2) \cdots (n-(k-1))}{k!} = \frac{n!}{k!(n-k)!}
\]
